\section{SAKI Design}
\label{sec:saki-design}

\begin{figure*}[!t]
  \centering
  \includegraphics[width=0.95\textwidth]{fig1_saki_framework.pdf}
  \caption{SAKI conserved rank-and-assign loop. Each window, the external
  controller scores tenants from previous-window demand and pressure signals,
  ranks them, assigns high/mid/low budget tiers, checks that the aggregate
  host budget is preserved, and pushes runtime rate-limit updates. The shown
  $4{:}8{:}4$ split at $11{:}6{:}1$\,MB/s is the main continuous setting.}
  \Description{Architecture diagram of SAKI. An external controller box at the
  top contains four pipeline steps: Score (rate each tenant from demand and
  pressure signals), Rank (sort by score), Assign budget tier (HIGH/MID/LOW
  at 11/6/1 MB/s), and Check and actuate (verify total budget conserved, push
  per-tenant rate limit). A closed loop runs every control window:
  observations flow up from the tenant pool to the controller, and budget
  updates flow down via the runtime rate limiter. The tenant pool at the
  bottom is one outer frame containing three tier cards, each with a labeled
  dashed sub-frame holding mini cylindrical RocksDB engine icons: four engines
  in the HIGH-budget tier, eight in the MID-budget tier, and four in the
  LOW-budget tier.}
  \label{fig:saki-framework}
\end{figure*}

SAKI is a host-level closed-loop controller that implements a conserved
rank-and-assign framework, sketched in Figure~\ref{fig:saki-framework}. The
framework is the method contribution: it preserves a fixed-budget invariant
and assignment template while allowing different observable score instances
to rank tenants. SAKI sits outside RocksDB, observes each tenant over window
\(w\), assigns budget tiers for window \(w+1\), and validates that the
aggregate budget is unchanged. The one-window delay is intentional: signals
from window \(w\) must remain informative for window \(w+1\).

\subsection{Conserved Rank-and-Assign Framework}

The framework decouples a budget-conserving assignment template from the score
instance used to rank tenants. The control loop, illustrated by the four
pipeline steps in Figure~\ref{fig:saki-framework}, is intentionally compact:
\[
\begin{array}{ll}
\textbf{Input:} & \{x_i(w)\}_{i=1}^{N},\; B,\; b_H,b_M,b_L,\; H,L \\
\textbf{Score:} & s_i(w) \leftarrow \mathrm{score}(x_i(w)) \\
\textbf{Rank:} & \text{sort tenants by descending } s_i(w) \\
\textbf{Assign:} & \text{top }H\rightarrow b_H,\; \text{bottom }L\rightarrow b_L,\;
\text{others}\rightarrow b_M \\
\textbf{Check:} & \sum_i b_i(w+1)=B \\
\textbf{Actuate:} & \text{apply } b_i(w+1) \text{ for the next window.}
\end{array}
\]
The tiered assignment makes the policy auditable. We distinguish
\emph{high-pressure} tenants, which the workload currently stresses, from
\emph{high-budget} tenants, which SAKI assigns to \(b_H\). A tier configuration
is budget-conserving when \(H(b_H-b_M)=L(b_M-b_L)\); the continuous main
setting uses \(4(11-6)=4(6-1)\), yielding the same 96\,MB/s aggregate budget as
sixteen all-mid tenants. Other tier counts or continuous allocations can use
the same rule; we fix auditable tierings to isolate the placement mechanism
within each experiment family.

The score instance instantiates the framework; it is not the framework itself.
This separation lets the paper make a framework claim and then test a small
family of auditable instances under the same conserved assignment rule. The
reference instance below is intentionally simple so the evaluation attributes
gains to online placement under conservation, not to an opaque learner or extra
capacity.

\subsection{Observations and Labels}

The observation vector is tenant-local, but the decision is host-wide. In the
continuous harness, \(x_i(w)\) includes offered write/read rates, completed
operations, write latency percentiles, compaction and flush bytes from an event
listener, estimated pending compaction bytes, L0 files, memtable bytes,
delayed-write rate, write-stop status, and the current assignment. These are
runtime measurements rather than workload labels.

The experiment generator retains each tenant's synthetic workload tier so the
evaluation can compute high-pressure benefit and high-set overlap. The
controller does not read that tier, a future phase identifier, or the next
high-pressure set. Offered write demand is an observed admission-side signal
from the previous window, rather than a prediction.

\subsection{Reference Score Instance and Family}

SAKI supports several score instances for ablation studies. The main continuous
policy uses a demand-aware reference instance. Let \(D_i\) be offered write
QPS, \(C_i\) completed write QPS, \(G_i=\max(0,D_i-C_i)\) the completion gap,
\(B_i\) the compaction output in MB, \(P_i\) pending compaction bytes in MB,
and \(L_i\) the number of L0 files capped at 6. The bounded tail term is
\[
T_i=\min(\mathrm{P99}_i/1000,80)+\min(\mathrm{P999}_i/5000,40),
\]
where latency is measured in microseconds. In demand mode, SAKI scores tenant
\(i\) as
\[
s_i =
6.00 D_i/1000
+0.60 G_i/1000
+0.05 T_i
+0.02 B_i
+0.02 P_i
+0.30 L_i .
\]
The coefficients reflect a design choice, not a universal optimum. Offered
write demand anchors active high-pressure tenants, while bounded residual
terms---completion gap, tail latency, compaction output, pending bytes, and L0
pressure---refine near-boundary tenants. This division targets two failures:
pressure-only scores can chase stale debt after demand moves, while demand-only
scores lack a maintenance signal for tenants that have just left a high-demand
phase. Residuals should not overturn a clearly separated demand signal, but
can stabilize rankings when margins are small. The evaluation tests this with
offline decision replay and online coefficient perturbations; in the main
continuous workload, replaying all 35 post-warmup decisions shows that the
demand term alone reproduces the full reference score's top-four assignment.

The score is fixed and auditable, but SAKI remains online: observations,
ranking, and assignment are recomputed every window. Keeping the rule fixed
separates the placement question from learning coefficients under delayed,
partially counterfactual feedback. The relevant design space is a family of
demand-aware instances whose primary term is offered write demand and whose
secondary pressure terms are bounded refinements. Section~\ref{sec:evaluation}
therefore tests demand-only replay, online anchor-scaled, residual-scaled, and
anchor-only instances, along with pressure-only negative controls, under the
same rank-and-assign template. This design yields two falsifiable predictions:
demand-aware instances should track moving pressure with high overlap and low
lag, while pressure-only instances should expose stale-debt lag. The
coefficients are not portable constants; deployments could add normalization or
policy layers while preserving budget conservation and demand observability.

The epoch harness uses a coarser, pressure-oriented score because each epoch
exposes clearer symptoms. The continuous harness is stricter: decisions are made
every short window while tenants stay active. Demand mode is the main
continuous policy; pressure and hybrid modes serve as ablations.

\subsection{Actuation and Stability}

The actuation path matches the harness. Epoch-level runs can change
background-job and rate-limit settings between epochs. Continuous runs keep
each tenant active and adjust only the runtime rate-limiter budget; \maxbgjobs
remains fixed at process launch. Section~\ref{sec:implementation} gives the
specific runtime API.

SAKI's framework-level stability rests on budget conservation, bounded score
terms, tiered assignment, a low-tier floor, and one-window updates. We evaluate
this envelope through high-set overlap, adaptation lag, total throughput,
bytes/write, and tier-level collateral effects. Production deployments could
add hysteresis, cooldowns, minimum-dwell rules, or SLO-weighted floors.
Per window, the controller aggregates \(N\) tenant observations, sorts tenants
by score, validates the conserved budget, and issues \(N\) budget updates. The
control-plane work is \(O(N\log N)\) and stays outside the foreground write
path. The paper does not include an overhead-scaling sweep at larger tenant
counts.

\subsection{Analytical Motivation}
\label{sec:analysis}

The design choices in the preceding subsections are motivated by a simple
analytical model. The model is not an optimizer for the coefficients; instead,
it isolates the tracking condition that the score family must satisfy and
explains the experimental divergence between demand-mode and pressure-mode
controllers reported in Section~\ref{sec:evaluation}.

\emph{Model.}
Let tenant \(i\) have offered write demand \(d_i(w)\), completed write rate
\(c_i(b)\) under budget \(b\), and normalized compaction debt \(q_i(w)\) in
equivalent writes. Budget tiers satisfy \(b_H>b_M>b_L\) with conservation
\(H b_H+(N-H-L)b_M+L b_L=B\). Backlog evolves as
\[
q_i(w+1)=\max\!\bigl(0,\; q_i(w)+T\,(d_i(w)-c_i(b_i(w)))\bigr),
\]
where \(T\) is the window length. We call tenant \(i\) \emph{high-pressure}
when \(d_i(w)\) exceeds the service rate achievable under \(b_M\); the deficit
under \(b_L\) is larger still.

\emph{Why static allocation can fail.}
Under equal allocation \(b_i=B/N\), if a high-pressure tenant satisfies
\(d_i>c_i(B/N)\), then \(q_i\) grows without bound until a write stall or
SLO violation occurs. This is the baseline problem: equal sharing allocates the
fixed budget without regard to which tenants are currently creating debt.

\emph{Proposition 1 (demand-aware one-window tracking).}
Consider a score instance in the rank-and-assign framework whose ordering can
be written as \(s_i=\alpha D_i(w)+\rho_i(w)\), where \(D_i(w)\) is the
observed offered demand and the residual term \(\rho_i(w)\) is computed from
bounded pressure and maintenance signals. Suppose the high-pressure set to be
served has at most \(H\) tenants and remains active into the next control
window. If, in window \(w\), every such tenant has offered demand at least
\(\Delta\) above every non-high tenant and
\[
\alpha\Delta > \max_{i,j}|\rho_i(w)-\rho_j(w)|,
\]
then the framework places those tenants in the high-budget tier in window
\(w+1\).

The proof is immediate from the ordering condition: for any high-pressure
tenant \(i\) and non-high tenant \(j\), \(s_i-s_j>0\), so the top-\(H\) rank
contains the high-pressure set. The one-window delay is unavoidable because the
controller only observes \(D_i(w)\) after window \(w\) completes. The condition
is sufficient, not universal: in the continuous main trials, offline replay
shows that the demand term alone produces the same top-four assignment as the
reference score in all 35 post-warmup decisions.

\emph{Proposition 2 (stale-debt lag for pressure-only ranking).}
Consider a monotone pressure-only instance that ranks tenants by backlog or
maintenance signals but does not observe offered demand. Immediately after a
workload drift, the instance cannot distinguish an exiting tenant with
historical debt from an entering tenant whose current demand is forming fresh
debt. If the exiting tenant starts the post-drift period with residual pressure
equivalent to \(q\) writes above the entrant, and the extra drain rate gained
by removing its high-budget slot is at most \(\Delta c=c(b_M)-c(b_L)\), then
the stale high-budget assignment can persist for
\(\Omega(q/(T\Delta c))\) windows under a monotone backlog-priority rule.

This is not a convergence theorem for all controllers; it is a structural
ambiguity in the pressure-only family. It predicts the measured divergence
between score instances: demand mode tracks with overlap \(3.00/4\) and zero
max lag, while pressure-only mode has overlap \(1.86/4\) and max lag 5
windows. The coefficients remain a deployment choice, but the
demand-observability condition is the mechanism that separates tracking from
stale-debt chasing.
